\chapter{HRBR Deterministic Scheduler: Lanes, Budgets, and Aging}

Signals answer \emph{what} should re-run. A scheduler answers \emph{when} and in what order.

This chapter is the bridge from ``reactivity is correct'' to ``reactivity feels fast''.
Correctness is table stakes.
The hard part in UI systems is correctness \emph{under time pressure}: you get 16ms for a 60Hz frame, and you don't control when the browser will deliver input.

If you're coming from React, the intuition you already have is that work can be prioritized and time-sliced.
React's documentation explains the conceptual split between render and commit (and why keeping commits responsive matters):
\href{https://react.dev/learn/render-and-commit}{https://react.dev/learn/render-and-commit}.

If you're coming from Solid, you already know the ``signals + effects'' world; this scheduler is the HRBR answer to: ``when should effects flush, and how do we avoid long synchronous waves?''
Background reading: \href{https://www.solidjs.com/docs/latest#scheduling}{https://www.solidjs.com/docs/latest\#scheduling}.

In HRBR, scheduling is part of the public contract. We want:

\begin{itemize}
  \item deterministic ordering (same inputs, same execution order)
  \item configurable latency via lanes (input vs transition vs idle)
  \item budget enforcement (avoid long flushes that block the main thread)
  \item starvation prevention (idle work should still eventually run)
\end{itemize}

This chapter is grounded in:

\begin{itemize}
  \item Implementation: \texttt{runtime/scheduler.ts}
  \item Specs: \texttt{runtime/\_\_tests\_\_/scheduler.test.ts}
  \item Devtools counters/events: \texttt{runtime/devtools.ts} (used by scheduler and signals)
\end{itemize}

\section{The contract: schedules tasks, flushes deterministically}

\subsection{Why a scheduler exists (not just ``because React has one'')}

Chapter~12 showed how signals can precisely compute \emph{what} is dirty.
However, once you have a list of dirty computations, you still need policy:

\begin{itemize}
  \item \textbf{Latency policy}: should an update run now (input), soon (default), or later (idle)?
  \item \textbf{Budget policy}: how long can the runtime keep running tasks before yielding?
  \item \textbf{Determinism policy}: when multiple tasks are eligible, what is the stable tie-breaker?
\end{itemize}

The HRBR scheduler makes these policies explicit and testable.

The scheduler exports three main concepts (see \texttt{runtime/scheduler.ts}):

\begin{itemize}
  \item \texttt{Lane}: \texttt{'sync' | 'input' | 'default' | 'transition' | 'idle'}
  \item \texttt{ScheduledTask}: \texttt{\{ id, lane, priority, timestamp, deadline, budgetMs, run \}}
  \item a factory: \texttt{createScheduler(options?) -> scheduler}
\end{itemize}

The returned scheduler exposes:

\begin{itemize}
  \item \texttt{schedule(lane, run, budgetOrOptions?)}
  \item \texttt{flush(\{ budgetMs \})}
  \item \texttt{hasPending()}
  \item \texttt{withBudget(budgetMs, fn)}
  \item \texttt{setFrameBudget(ms)} and \texttt{setDefaultBudget(ms)}
\end{itemize}

The tests treat these as the source of truth.

\subsection{A concrete mental model: 5 queues + 1 deterministic pick loop}

At runtime you can picture five arrays (one per lane), and a function that repeatedly:

\begin{enumerate}
  \item optionally promotes one starving head task
  \item picks the next lane in fixed order
  \item shifts one task and runs it
  \item stops when its time budget is exhausted
\end{enumerate}

Even the ``browser integration'' is just a \texttt{requestFlush(cb)} option (\texttt{setTimeout}, \texttt{MessageChannel}, or \texttt{requestAnimationFrame}).

\begin{lstlisting}[language=JavaScript]
import { createBrowserScheduler } from './runtime/scheduler'

// In the browser you can pick a flushing strategy.
const scheduler = createBrowserScheduler({
  strategy: 'messageChannel',
  defaultBudgetMs: 5,
})

scheduler.schedule('input', () => {
  // keep interactions snappy
})

scheduler.schedule('idle', () => {
  // background bookkeeping
})
\end{lstlisting}

\paragraph{A tiny scheduler contract}

We can summarize what the runtime scheduler promises in four bullets:

\begin{itemize}
  \item \textbf{Pick order is deterministic}: same task creation order + same time evolution \textrightarrow same execution order.
  \item \textbf{Latency is explicit}: lane names are the only priority vocabulary.
  \item \textbf{Work is time-sliced}: every flush loop has a budget, and work may continue later.
  \item \textbf{Starvation is mitigated}: low-priority tasks can be promoted based on age/deadline.
\end{itemize}

\section{Lanes: latency categories and fixed priority order}

Lane priority is encoded as a small integer map:

\begin{verbatim}
const LANE_PRIORITY = {
  sync: 0,
  input: 1,
  default: 2,
  transition: 3,
  idle: 4,
}
\end{verbatim}

But the scheduler's actual selection order is hard-coded as an array:

\begin{verbatim}
['sync', 'input', 'default', 'transition', 'idle']
\end{verbatim}

Spec: \texttt{runs tasks by lane priority (sync > input > default > transition > idle)}.

\subsection{What lanes are (and what they are not)}

It's easy to overthink lanes.
In HRBR they're not threads, not preemption, and not ``multiprocessing''.
A lane is just a \emph{label} on a task that determines its relative ordering.

This matters because it keeps the model honest:

\begin{itemize}
  \item tasks still run to completion
  \item the only question is: which task runs \emph{next} and when do we continue later?
\end{itemize}

\subsection*{Diagram: lane priority at a glance}

\begin{center}
\begin{tikzpicture}[
  node distance=7mm,
  box/.style={draw,rounded corners,fill=RelaxLight,inner sep=5pt,align=center},
  arrow/.style={-Latex,thick}
]
  \node[box] (sync) {sync\\(run now)};
  \node[box, below=of sync] (input) {input\\(user interaction)};
  \node[box, below=of input] (def) {default\\(normal updates)};
  \node[box, below=of def] (trans) {transition\\(deferred UI)};
  \node[box, below=of trans] (idle) {idle\\(background)};

  \draw[arrow] (sync) -- (input);
  \draw[arrow] (input) -- (def);
  \draw[arrow] (def) -- (trans);
  \draw[arrow] (trans) -- (idle);
\end{tikzpicture}
\end{center}

This ordering applies unless aging/deadline promotion temporarily pushes work upward.

\subsection{What is a lane in practice?}

Lanes are not preemption.
They’re an ordering and continuation policy:

\begin{itemize}
  \item tasks run to completion one-by-one
  \item lane selection chooses which queue head runs next
  \item aging/deadlines can temporarily override lane order via promotion
\end{itemize}

\subsection{The sync lane}

A \texttt{sync} task is special: \texttt{schedule('sync', ...)} triggers an immediate flush.

In \texttt{schedule()}:

\begin{itemize}
  \item the task is still enqueued
  \item then the scheduler calls \texttt{flush(\{ budgetMs: Infinity \})}
\end{itemize}

This gives a "run now" lane while still reusing the same task machinery.

\paragraph{Why still enqueue sync tasks?}

Because determinism still matters.
Enqueueing before flushing means sync work still participates in the same stable ordering rules if multiple sync tasks are scheduled.

\section{Flush strategies: how work gets a time slice}

The scheduler accepts \texttt{requestFlush(cb)} in its options.
If you don't provide one, it uses \texttt{createRequestFlush('timeout')}.

\subsection{createRequestFlush(strategy)}

\texttt{createRequestFlush()} supports:

\begin{itemize}
  \item \texttt{'timeout'} (universal fallback): \texttt{setTimeout(cb, 0)}
  \item \texttt{'messageChannel'}: uses \texttt{MessageChannel} if available
  \item \texttt{'raf'}: uses \texttt{requestAnimationFrame} if available
\end{itemize}

Specs in \texttt{runtime/\_\_tests\_\_/scheduler.test.ts} cover:

\begin{itemize}
  \item \texttt{createRequestFlush(messageChannel) uses MessageChannel when available}
  \item \texttt{createRequestFlush(raf) uses requestAnimationFrame when available}
\end{itemize}

\subsection{createBrowserScheduler() convenience}

\texttt{createBrowserScheduler({ strategy })} is a wrapper that wires the chosen strategy into \texttt{requestFlush}.

Spec: \texttt{createBrowserScheduler wires a strategy into requestFlush (smoke)}.

\paragraph{Why these strategies exist}

The scheduler is designed to be embeddable:

\begin{itemize}
  \item in Node-like environments (\texttt{setTimeout})
  \item in browsers where \texttt{MessageChannel} is often lower-latency than timeouts
  \item in UIs where \texttt{requestAnimationFrame} aligns work with the next frame
\end{itemize}

\section{Deterministic ordering inside a lane}

Within a lane, tasks are sorted with explicit tie-breakers:

\begin{enumerate}
  \item earlier \texttt{deadline}
  \item earlier \texttt{timestamp}
  \item lower \texttt{id}
\end{enumerate}

This ordering is implemented by \texttt{sortQueueInLane()}:

\begin{verbatim}
q.sort((a, b) => (a.deadline - b.deadline)
  || (a.timestamp - b.timestamp)
  || (a.id - b.id))
\end{verbatim}

Spec: \texttt{orders tasks within a lane by explicit deadline (then timestamp/id)}.

\subsection{Where deadlines come from}

If you don't provide an explicit \texttt{deadline}, \texttt{schedule()} computes:

\begin{verbatim}
deadline = timestamp + budgetMs
\end{verbatim}

So tasks with smaller budgets naturally get earlier deadlines and bubble forward within their lane.
This is a useful ``do small things first'' heuristic that is still deterministic.

\subsection{Three tie-breakers, one guarantee (and why we need all three)}

It’s tempting to think only deadlines matter, but the scheduler needs a \emph{total} ordering to avoid accidental nondeterminism.
Inside a single lane, two tasks can easily share a deadline (same budget) and timestamp (scheduled in the same tick).
The created \texttt{id} is the final ``stable sort'' key.

The guarantee we want is:

\begin{quote}
If you schedule the same tasks in the same order, you will flush them in the same order.
\end{quote}

This is why the implementation stores \texttt{id} and sorts by \texttt{(deadline, timestamp, id)}.

\section{Budgets: splitting work over multiple flushes}

There are two different "budget" knobs:

\begin{itemize}
  \item \texttt{defaultBudgetMs}: the per-task budget used to compute default \texttt{deadline} and default \texttt{budgetMs}
  \item \texttt{frameBudgetMs}: a global cap for \emph{scheduled flush loops}
\end{itemize}

\subsection{What flush(budgetMs) does}

\texttt{flush(\{ budgetMs \})}:

\begin{itemize}
  \item emits devtools events \texttt{flushStart} and \texttt{flushEnd}
  \item runs tasks in priority order until:
    \begin{itemize}
      \item queues are empty, or
      \item elapsed time exceeds \texttt{budgetMs}
    \end{itemize}
  \item if budget is exceeded \emph{and} work remains, it schedules a continuation flush
\end{itemize}

Specs:

\begin{itemize}
  \item \texttt{budget enforcement: flush splits work across multiple flush calls when budget is tight}
  \item \texttt{setFrameBudget caps the budget used by scheduled flushes}
\end{itemize}

\subsection*{Diagram: the flush loop (budgeted work)}

\begin{center}
\begin{tikzpicture}[
  node distance=8mm,
  box/.style={draw,rounded corners,fill=RelaxLight,inner sep=5pt,align=left},
  arrow/.style={-Latex,thick}
]
  \node[box] (start) {flush(\{budgetMs\}) starts};
  \node[box, below=of start] (pick) {pick next task\\(lane priority + promotion)};
  \node[box, below=of pick] (run) {run task};
  \node[box, below=of run] (check) {elapsed > budget?};
  \node[box, below left=of check, xshift=14mm] (cont) {yes:\\schedule continuation flush\\work remains};
  \node[box, below right=of check, xshift=-14mm] (done) {no:\\continue until empty};

  \draw[arrow] (start) -- (pick);
  \draw[arrow] (pick) -- (run);
  \draw[arrow] (run) -- (check);
  \draw[arrow] (check) -- (cont);
  \draw[arrow] (check) -- (done);
\end{tikzpicture}
\end{center}

\subsection{setFrameBudget(ms)}

Even if tasks provide their own budgets, scheduled flushes are capped by \texttt{frameBudgetMs}.

The spec demonstrates this by setting \texttt{frameBudgetMs = 1} and scheduling tasks that each consume \texttt{2ms}; only one task runs per flush, and the second flush runs the remainder.

\subsection{Frame budget versus task budget (important distinction)}

In \texttt{runtime/scheduler.ts}, \texttt{budgetMs} is stored on each task, but it isn't directly used by \texttt{flush()}.
Instead:

\begin{itemize}
  \item \textbf{task budget} influences default deadlines (and therefore ordering)
  \item \textbf{flush budget} (\texttt{flush(\{budgetMs\})} or \texttt{frameBudgetMs}) limits how long the flush loop runs before continuation is scheduled
\end{itemize}

The test suite focuses on flush budgets because those directly affect responsiveness.

\subsection{Walkthrough: tight budget slicing (mapping a test to code)}

Read the spec named \texttt{budget enforcement: flush splits work across multiple flush calls when budget is tight}.
It schedules three tasks; each task increments a manual clock by 3ms.
Then it runs two flush calls with \texttt{budgetMs = 5}.

The implementation in \texttt{runtime/scheduler.ts} does exactly this:

\begin{enumerate}
  \item stores \texttt{start = now()}
  \item runs one task at a time
  \item recomputes \texttt{elapsed = now() - start}
  \item stops once \texttt{elapsed >= budgetMs}
\end{enumerate}

So with 3ms per task:

\begin{itemize}
  \item task 1 leaves us at 3ms (still under 5ms)
  \item task 2 leaves us at 6ms (over budget) \textrightarrow stop and continue later
  \item task 3 remains queued for the next flush
\end{itemize}

This is "cooperative scheduling": tasks cooperate by being small and letting the runtime yield.
If you want more intuition about why we care, the web performance RAIL model is a good framing for latency budgets:
\href{https://web.dev/rail/}{https://web.dev/rail/}.

\subsection{Timeline example: lane priority, aging, and budget all at once}

The trickiest behavior is when \emph{multiple policies apply}: lane order, promotion, and a flush budget.
Let's walk a concrete timeline using the same concepts as \texttt{makeManualScheduler()} in the tests.

Assume:

\begin{itemize}
  \item \texttt{AGING\_MS = 250}
  \item \texttt{frameBudgetMs = 5}
  \item each task costs \texttt{3ms} (it advances the manual clock)
\end{itemize}

We schedule:

\begin{lstlisting}[language=JavaScript]
// T = 0
scheduler.schedule('idle', () => tick(3))

// time passes; the idle head is now "old"
tick(300)

// scheduling default work ensures a flush is scheduled
scheduler.schedule('default', () => tick(3))
\end{lstlisting}

Now what happens during \texttt{flush()}?

\begin{enumerate}
  \item \textbf{Promotion step}: before lane selection, the scheduler checks the idle head.
        It has aged, so it is promoted to \texttt{transition}.
  \item \textbf{Lane selection}: the scheduler picks the highest-priority non-empty lane.
        At this point we have \texttt{default} (new) and \texttt{transition} (promoted).
        The lane order says \texttt{default} wins.
  \item \textbf{But why do tests show idle can run first?} Because the implementation promotes by \emph{moving the task to a higher-priority lane} before selection.
        Across repeated picks (and depending on the exact mix in queues), promotion can change which lane is "non-empty" at the top.
\end{enumerate}

The important lesson isn't the specific interleaving---it's that the promotion rule is deterministic and testable.
When you change the policy, the tests will immediately surface the new ordering.

\subsection{Edge cases to keep in mind}

When you build APIs on top of this scheduler (signals, blocks), three edge cases matter:

\begin{itemize}
  \item \textbf{Re-entrancy}: a running task may schedule more tasks. Determinism requires those tasks get consistent timestamps/ids.
  \item \textbf{Budget 0}: \texttt{setFrameBudget(0)} makes scheduled flushes yield aggressively. That's allowed, but it can create many continuation flushes.
  \item \textbf{Clock source}: using \texttt{performance.now()} in production versus a manual clock in tests is fine, but it means ``timestamp equality'' is common only in tests.
\end{itemize}

\section{Deadlines and aging: starvation prevention via promotion}

HRBR's scheduler includes an aging/promotion mechanism.

In the implementation:

\begin{itemize}
  \item \texttt{AGING\_MS = 250}
  \item lanes eligible for promotion checks: \texttt{idle}, \texttt{transition}, \texttt{default}
  \item promotion moves one step toward higher priority:
    \texttt{idle -> transition -> default -> input}
\end{itemize}

Promotion runs \emph{before} picking the next task, and it only considers the \textbf{queue head} in each lane.

\subsection*{Diagram: promotion ladder}

\begin{center}
\begin{tikzpicture}[
  node distance=7mm,
  box/.style={draw,rounded corners,fill=RelaxLight,inner sep=5pt,align=center},
  arrow/.style={-Latex,thick}
]
  \node[box] (idle) {idle};
  \node[box, right=of idle] (trans) {transition};
  \node[box, right=of trans] (def) {default};
  \node[box, right=of def] (input) {input};
  \draw[arrow] (idle) -- (trans);
  \draw[arrow] (trans) -- (def);
  \draw[arrow] (def) -- (input);
\end{tikzpicture}
\end{center}

The scheduler promotes a task when it's overdue (deadline reached) or old (aging threshold). This is the mechanism behind the starvation-prevention specs.

A task is promoted if:

\begin{itemize}
  \item it is overdue: \texttt{now() >= head.deadline}, or
  \item it is old: \texttt{now() - head.timestamp > AGING\_MS}
\end{itemize}

Specs:

\begin{itemize}
  \item \texttt{ages/pushes starving tasks upwards (idle -> transition)}
  \item \texttt{promotes a task when its deadline is reached}
  \item \texttt{starvation prevention: repeated promotions eventually let idle work run even under continuous input load}
\end{itemize}

One subtlety the tests call out:

\begin{quote}
Promotion happens before lane selection, so an aged/overdue idle task can execute before later-scheduled default work.
\end{quote}

This is a deliberate fairness trade-off.

\subsection{Fairness without nondeterminism}

The promotion rule checks only \textbf{queue heads} and promotes at most one task per candidate lane per pick.
That has two benefits:

\begin{itemize}
  \item it's cheap (no scanning the entire queue)
  \item it stays deterministic (the head is a deterministic choice)
\end{itemize}

\subsection{Why promotion checks only queue heads}

The promotion policy is intentionally restricted to heads:

\begin{itemize}
  \item It's \textbf{predictable}: given the same inputs and the same clock, you'll promote the same tasks.
  \item It's \textbf{O(1)} per lane: no scanning and no ``who is most starving'' heuristics.
  \item It preserves within-lane ordering as much as possible.
\end{itemize}

This is a classic engineering trade-off: accept a slightly simplified fairness model to buy determinism and low overhead.

\section{withBudget(): bridging user work and scheduler work}

Real apps do work outside of the scheduler too (rendering, event handlers, parsing).
\texttt{withBudget(budgetMs, fn)} is a helper that:

\begin{itemize}
  \item runs \texttt{fn}
  \item if \texttt{fn} exceeds the budget and tasks remain, it schedules a continuation flush
\end{itemize}

Spec: \texttt{withBudget schedules continuation when user code exhausts budget and tasks remain}.

This is the "be nice to the main thread" glue between imperative work and queued reactive work.

\subsection{A practical use case: big event handler + pending effects}

Imagine an event handler that does some synchronous work (parsing, filtering, sorting) and also triggers signal writes.
Those writes schedule computations, but you still want control: if your handler took too long, queue a flush continuation.

\begin{lstlisting}[language=JavaScript]
import { withBudget } from './runtime'

withBudget(4, () => {
  // expensive user work
  // ... parse, compute, validate ...
  // signals may schedule effects while we do this
})
// if the 4ms budget was exceeded and work is pending,
// the scheduler will ensure another flush is scheduled.
\end{lstlisting}

\subsection{Global helpers in \texttt{runtime/index.ts}}

The runtime also exports shared budgeting helpers that use a module-level global scheduler:

\begin{itemize}
  \item \texttt{withBudget(budgetMs, fn)}
  \item \texttt{setFrameBudget(ms)}
\end{itemize}

This is covered by the smoke test \texttt{runtime withBudget() helper returns the callback result and is callable}.

\section{Devtools and instrumentation hooks}

The scheduler calls \texttt{emitDevtoolsEvent()} (\texttt{runtime/devtools.ts}) when flushing:

\begin{itemize}
  \item \texttt{flushStart}
  \item \texttt{flushEnd\{ durationMs \}}
\end{itemize}

These feed counters like \texttt{flushes} and \texttt{flushTimeMs}.

Signals also emit \texttt{scheduleComputation} events when a computation is scheduled.

(We’ll use these counters later when measuring block update costs.)

\subsection{Why ``events + counters'' is a good runtime pattern}

Two competing goals exist:

\begin{itemize}
  \item keep the runtime fast (no heavy logging)
  \item make the runtime observable (so you can measure and debug)
\end{itemize}

The implementation in \texttt{runtime/devtools.ts} uses a tiny event stream plus aggregated counters.
That means you can do ``always-on'' cheap counting (flush count/time) and keep deeper instrumentation behind a toggle.

\section{Walking the tests: what they prove}

The runtime scheduler suite is unusually readable because it uses a manual clock and a manual flush queue.
That makes time and continuation behavior deterministic.

\subsection{Core semantics (\texttt{runtime/\_\_tests\_\_/scheduler.test.ts})}

\begin{description}
  \item[runs tasks by lane priority] Proves the lane order and that \texttt{sync} executes immediately.
  \item[orders tasks within a lane by explicit deadline] Proves the sort tie-breakers inside a lane.
  \item[respects budget and continues in a later flush] Proves a scheduled flush loop may stop early and reschedule continuation.
  \item[budget enforcement: flush splits work] Proves deterministic time slicing under a small budget.
  \item[setFrameBudget caps scheduled flush budget] Proves continuation flushes use \texttt{frameBudgetMs} as a cap.
  \item[ages/pushes starving tasks upwards] Proves age-based promotion and the fairness trade-off (promoted work can run before later default work).
  \item[promotes a task when its deadline is reached] Proves deadline-based promotion independent of age.
  \item[starvation prevention under input load] Proves repeated promotions eventually allow idle work to run.
  \item[createRequestFlush(messageChannel/raf)] Proves strategy selection uses platform APIs when present.
  \item[createBrowserScheduler smoke] Proves wiring strategy into \texttt{requestFlush}.
\end{description}

\subsection{How to read these tests like a spec}

The tests are valuable because they remove nondeterminism from time.
The helper \texttt{makeManualScheduler()} provides:

\begin{itemize}
  \item a manual clock: \texttt{tick(ms)}
  \item a list of pending flush callbacks: \texttt{runNextFlush()} / \texttt{runAllFlushes()}
\end{itemize}

So when you see an assertion like \texttt{expect(order).toEqual([...])}, you're basically reading the scheduler's language definition.

\section{Online resources}

\begin{itemize}
  \item MDN: \href{https://developer.mozilla.org/en-US/docs/Web/API/MessageChannel}{\texttt{MessageChannel}}
  \item MDN: \href{https://developer.mozilla.org/en-US/docs/Web/API/window/requestAnimationFrame}{\texttt{requestAnimationFrame}}
  \item MDN: \href{https://developer.mozilla.org/en-US/docs/Web/JavaScript/Event_loop}{Event loop (tasks, microtasks)}
  \item React (conceptual parallel): \href{https://react.dev/learn/render-and-commit}{Render/commit (scheduling intuition)}
  \item Cooperative background work (conceptual): \href{https://developers.google.com/web/updates/2015/08/using-requestidlecallback}{\texttt{requestIdleCallback}}
\end{itemize}

\section{Takeaways}

\begin{itemize}
  \item Lanes provide a small, explicit latency model with a deterministic priority order.
  \item Each lane is internally sorted by deadline, then timestamp, then id.
  \item Budgeting is enforced at flush time to time-slice work.
  \item Aging/promotion makes starvation unlikely without making scheduling non-deterministic.
  \item \texttt{withBudget} lets user code cooperate with the scheduler.
\end{itemize}
